% Generated draft: requires mathematical and source review.
\documentclass[11pt]{article}
\usepackage[margin=25mm]{geometry}
\usepackage{fontspec}
\setmainfont{Latin Modern Roman}
\usepackage{amsmath,amssymb,mathtools,unicode-math}
\setmathfont{Latin Modern Math}
\usepackage{enumitem,needspace,xurl,hyperref}
\setlength{\parindent}{0pt}
\setlength{\parskip}{5pt}
\setlength{\emergencystretch}{3em}
\newcommand{\N}{\mathbb{N}}
\newcommand{\Z}{\mathbb{Z}}
\newcommand{\Q}{\mathbb{Q}}
\newcommand{\R}{\mathbb{R}}
\newcommand{\C}{\mathbb{C}}
\newcommand{\F}{\mathbb{F}}
\newcommand{\A}{\mathbb{A}}
\newcommand{\PP}{\mathbb{P}}
\newcommand{\E}{\mathbb{E}}
\newcommand{\eps}{\varepsilon}
\newcommand{\dd}{\,\mathrm{d}}
\begin{document}
\section*{UnsolvedMath ID: 2302030}
\textbf{Source problem code:} catalogue entry
\subsection{Definitions and mathematical statement}
For each positive integer $k$, let $S_k\subset\mathbb C$ have no finite limit point; equivalently, its intersection with every compact subset of $\mathbb C$ is finite. For an entire function $f$ and an integer $n\geq0$, $f^{(n)}$ denotes its $n$th complex derivative, with $f^{(0)}=f$.

\textbf{Question (P. Erd\H{o}s).} For every sequence $(S_k)_{k\geq1}$ of such sets, do there exist a sequence of derivative orders $(n_k)_{k\geq1}$ and an entire function $f$ such that
\[
\text{for every }k\geq1\text{ and every }z\in S_k,
\qquad f^{(n_k)}(z)=0?
\]
Only inclusion of $S_k$ in the zero set is required; no assertion that these are all the zeros or that they are simple is made.

The source does not restrict $f$ to be nonzero, nonconstant, or transcendental, and does not require distinct or increasing derivative orders. Those omissions permit vacuous polynomial or identically zero cases and must be resolved before treating the request as a nontrivial open interpolation problem. The conventional derivative orders above are nonnegative integers, but the source does not say whether $0$ is permitted. No missing nontriviality or growth condition is invented. The supplied summary reports the problem open as of the 2018 edition, with no recent resolution located.
\subsection{Short English statement}
Prescribe discrete sets on which selected derivatives of one entire function vanish. The literal statement lacks the nontriviality and derivative-order restrictions needed to identify the intended open problem.
\subsection{Sources}
\begin{itemize}
\item Hayman, \emph{Research Problems in Function Theory} (2018), Problem 2.30, attributed to P. Erd\H{o}s, \url{https://arxiv.org/abs/1809.07200}; source-TeX extraction, accessed 2026-07-29. The 2018 update was retained for status triage.
\item Local registry \texttt{data/unsolvedmath/problems.json}, permanent ID $2302030$, source code catalogue entry, \emph{Research Problems in Function Theory --- Problem 2.30}. The record labels this open and OPEN-TRIAGE, with NEEDS\_REVIEW source-status and rights qualifications. No separate result citation is supplied.
\end{itemize}
\end{document}
